\section{Analysis of WG-KV's Input-Dependent Admission Patterns}
\label{app:dynamic}

To demonstrate that WG-KV learns an input-dependent admission policy rather than degenerating into input-agnostic head specialization, we provide both qualitative and quantitative analyses in this section.

\subsection{Qualitative analysis}
\label{app:qualitative}

To show that WG-KV adapts its memory allocation based on the input context, \cref{fig:gate} visualizes the distribution of KV cache sizes across all attention heads using Llama-3.1-8B with $\lambda=0.08$. We compare two distinct long-context tasks: \textit{(a) Code Summarization}, utilizing a Python sample from The Stack \cite{thestack}, and \textit{(b) HTML to TSV}, employing a structured data extraction sample from LongProc \cite{longproc}.

\paragraph{Observations.} The heatmaps reveal distinct retention characteristics tailored to the input task:
\begin{enumerate}
    \item \textbf{Task sensitivity.} For \textit{(a) Code Summarization}, the cache distribution is relatively uniform, indicating that understanding code dependencies requires retaining context across diverse heads. In contrast, for \textit{(b) HTML to TSV}, the pattern is significantly sparser, suggesting that the relevant information for this task is highly concentrated. The model effectively filters out most tokens, likely preserving primarily the critical structural cues (e.g., tags) to perform the extraction.
    \item \textbf{Head-specific decisions.} Consistent with \cref{fig:distribution}, the admission policy is fine-grained. Specific heads retain nearly full history, while adjacent heads may discard most tokens.
\end{enumerate}

These findings demonstrate that WG-KV adaptively modulates its memory allocation in response to the semantic structures of the task, dynamically transitioning between high-retention and high-sparsity modes.

\subsection{Quantitative analysis}

To further justify that admission decisions are dynamically driven by the input rather than relying on input-agnostic head profiles, we conduct a cross-sequence replay experiment. We randomly sample 1,000 pairs of 8K-length sequences $(S_1,S_2)$ from the PG-19 \cite{pg19} long-document dataset and feed them into Llama-3.1-8B with $\lambda=0.08$. We measure the average MSE of the last-layer hidden states for $S_1$ against the standard full-attention model under three policies:

\begin{enumerate}
    \item \textbf{Own Policy:} Using $S_1$'s own dynamically generated admission decisions.
    \item \textbf{Replay Policy:} Forcing $S_1$ to adopt the exact admission decisions generated by $S_2$.
    \item \textbf{Random Policy:} Randomly admitting tokens for $S_1$ while matching the exact retention ratio of $S_2$.
\end{enumerate}

If WG-KV merely degenerated into learning input-agnostic head specialization, the \textit{Replay Policy} should yield an MSE close to \textit{Own Policy} and outperform \textit{Random Policy}. However, as shown in \cref{tab:replay}, the \textit{Replay Policy} performs almost identically to the \textit{Random Policy} and is significantly worse than \textit{Own Policy}. This stark contrast quantitatively confirms that WG-KV's admission decisions are strictly input-dependent and dynamically tailored to the context of the given sequence.

\tabReplay{t}{0.9}